% Sources: W01-W06 and L48-L53.
\boardsection{6}{充分统计量}
1．定义法。

2．因子分解定理。
\[
L(\theta\mid x)=h(x)g(T(x),\theta).
\]

\block{指数分布}
\[
f(x\mid\theta)=\theta e^{-\theta x}\ind_{\{x>0\}},\quad\theta>0,
\qquad T=\sum_{i=1}^nX_i.
\]
考虑$(X_1,\ldots,X_n)$到$(X_1,\ldots,X_{n-1},T)$的变换，
\[
X_n=T-\sum_{i=1}^{n-1}X_i,
\]
雅可比行列式的绝对值为1。
\begin{align*}
f(x_1,\ldots,x_{n-1},t\mid\theta)
&=\theta^n
\exp\left\{-\theta\left[\sum_{i=1}^{n-1}x_i+
\left(t-\sum_{i=1}^{n-1}x_i\right)\right]\right\}\\
&\quad\times\ind_{\{x_1>0,\ldots,x_{n-1}>0,\ t-\sum_{i=1}^{n-1}x_i>0\}}\\
&=\theta^ne^{-\theta t}
\ind_{\{x_1>0,\ldots,x_{n-1}>0,\ \sum_{i=1}^{n-1}x_i<t\}}.
\end{align*}
\[
T\sim\Ga(n,\theta),\qquad
f_T(t\mid\theta)=\frac{\theta^n}{\Gamma(n)}t^{n-1}e^{-\theta t},
\quad t>0,\quad\Gamma(n)=(n-1)!.
\]
\begin{align*}
&f(x_1,\ldots,x_{n-1}\mid T=t,\theta)\\
&=\frac{f(x_1,\ldots,x_{n-1},t\mid\theta)}{f_T(t\mid\theta)}\\
&=\frac{\theta^ne^{-\theta t}
\ind_{\{x_i>0,\ \sum_{i=1}^{n-1}x_i<t\}}}
{\frac{\theta^n}{(n-1)!}t^{n-1}e^{-\theta t}}\\
&=\frac{(n-1)!}{t^{n-1}}
\ind_{\{x_i>0,\ \sum_{i=1}^{n-1}x_i<t\}}.
\end{align*}
该条件分布与$\theta$无关，$T=\sum_iX_i$为充分统计量。

\block{例2 正态分布}
样本的联合密度为
\begin{align*}
L(\mu,\sigma^2\mid x)
&=\prod_{i=1}^n\frac1{\sqrt{2\pi\sigma^2}}
\exp\left\{-\frac{(x_i-\mu)^2}{2\sigma^2}\right\}\\
&=(2\pi\sigma^2)^{-n/2}
\exp\left\{-\frac1{2\sigma^2}\sum_i(x_i-\mu)^2\right\}.
\end{align*}
而
\[
\sum_i(x_i-\mu)^2=\underbrace{\sum_i(x_i-\bar x)^2}_{Q}
+n(\bar x-\mu)^2.
\]
故
\[
L(\mu,\sigma^2\mid x)
=(2\pi\sigma^2)^{-n/2}
\exp\left\{-\frac{Q+n(\bar x-\mu)^2}{2\sigma^2}\right\}
=g(T(x),\mu,\sigma^2)h(x),
\]
\[
h(x)=1.
\]
由因子分解定理，$T(x)=(\bar x,Q)$是$(\mu,\sigma^2)$的充分统计量。

\block{例1.6.2}
设$X_1,\ldots,X_n\overset{\mathrm{iid}}{\sim}N(\theta,1)$，
$\bar X=\frac1n\sum_iX_i$为$\theta$的充分统计量。
取先验$\theta\sim N(0,\tau^2)$。
\[
X_i\mid\theta\sim N(\theta,1),\qquad
\bar X\mid\theta\sim N(\theta,1/n),\qquad
\Var(\bar X\mid\theta)=1/n.
\]
\[
f(\bar x\mid\theta)=\sqrt{\frac n{2\pi}}
\exp\left\{-\frac n2(\bar x-\theta)^2\right\},\qquad
L(\theta\mid\bar x)\propto
\exp\left\{-\frac n2(\bar x-\theta)^2\right\}.
\]
\[
\pi(\theta)=\frac1{\sqrt{2\pi\tau^2}}
\exp\left\{-\frac{\theta^2}{2\tau^2}\right\}
\propto\exp\left\{-\frac{\theta^2}{2\tau^2}\right\}.
\]
代入似然与先验核，有
\begin{align*}
\pi(\theta\mid\bar x)
&\propto L(\theta\mid\bar x)\pi(\theta)\\
&\propto\exp\left\{-\frac12
\left[n(\bar x-\theta)^2+\frac{\theta^2}{\tau^2}\right]\right\}\\
&=\exp\left\{-\frac12
\left[n\bar x^2-2n\bar x\theta+n\theta^2+\tau^{-2}\theta^2\right]\right\}\\
&\propto\exp\left\{-\frac12[(n+\tau^{-2})\theta^2-2n\bar x\theta]\right\}.
\end{align*}
对其配方：
\begin{align*}
(n+\tau^{-2})\theta^2-2n\bar x\theta
&=(n+\tau^{-2})\left[\theta^2-\frac{2n\bar x}{n+\tau^{-2}}\theta\right]\\
&=(n+\tau^{-2})\left(\theta-\frac{n\bar x}{n+\tau^{-2}}\right)^2
-\frac{n^2\bar x^2}{n+\tau^{-2}}.
\end{align*}
\[
\pi(\theta\mid\bar x)\propto
\exp\left\{-\frac{n+\tau^{-2}}2
\left(\theta-\frac{n\bar x}{n+\tau^{-2}}\right)^2\right\}.
\]
\[
\theta\mid\bar x\sim
N\left(\frac{n\bar x}{n+\tau^{-2}},\frac1{n+\tau^{-2}}\right)
=N\left(\frac{n\tau^2}{n\tau^2+1}\bar x,\frac{\tau^2}{n\tau^2+1}\right).
\]
\[
\E(\theta\mid\bar x)=
\frac n{n+\tau^{-2}}\bar x+\frac{\tau^{-2}}{n+\tau^{-2}}\times0.
\]
